% !TEX root = ../main.tex

\appendix
\section{Missing Runs and the Choice of Classifier}
\label{sec:appendix}

\paragraph{Missing runs.} Of the 195 planned benchmark runs per configuration and classifier, 69 to 87\% finished for the OBL start, 85 to 100\% for the sparse start, and 68 to 82\% for NSGAII-MIIP (Supplementary Table~S4); all 80 runs of the designed experiment finished. Under the decision tree, 47 benchmark runs reached the six-hour limit (6 OBL, 30 sparse, and 11 NSGAII-MIIP runs), under LightGBM 64 (36, 1, and 27), and under the random forest none. The combinations that were not run remove 25 OBL and 35 NSGAII-MIIP runs from each group. Missing runs are thus concentrated in the configurations and on the wide datasets where reach is low. Figure~\ref{fig:missing} repeats the main comparison under the three treatments of Section~\ref{sec:statistics}. In every group and under every treatment, the OBL start never reaches the limit in more seeds than the sparse start, and the absolute reach share $A$ is at least 97.9\%. These treatments are scenarios; they do not rule out that runs that did not finish would have changed the comparison in another way. Counting every missing run as reaching nothing moves the mean difference in reach from 0.68, 0.64, and 0.54 to 0.59, 0.68, and 0.67. One decision outside the main comparison changes: with only complete datasets, the AUF difference between NSGAII-MIIP and the opposition-based start under LightGBM has Holm $p = 0.055$ instead of $0.029$.

\begin{figure}[!htbp]
  \centering
  \includegraphics{fig_missing}
  \caption{The main comparison persists under the tested treatments of missing runs. (a) Mean difference in reach, sparse minus OBL start; (b) the absolute share $A$ of the AUF difference due to reach. As registered: finished runs only; complete datasets: datasets on which all seeds finished; missing as no reach: missing runs count as reaching nothing; shared datasets: datasets on which every configuration has runs}
  \label{fig:missing}
\end{figure}

\paragraph{The classifier.} On the 26 datasets on which both starts have runs under all three classifiers, the AUF lead of the sparse start is 0.465 under the decision tree, 0.462 under the random forest, and 0.448 under LightGBM (Supplementary Fig.~S1). The Friedman test gives $p = 0.007$, but only the difference between the tree and the forest remains after Holm correction ($+0.003$, $p = 0.022$), and the absolute reach share $A$ is 99.7\%, 99.4\%, and 99.5\%. The share of opposition-based runs that reach ten features on its 27 common datasets is 0.36, 0.42, and 0.44 (Friedman $p = 0.082$). The classifier changes the scores more than the comparison: the sparse start's own AUF is 0.04 lower under the tree than under the other two (Holm $p < 0.001$).
